% Build in this directory: pdflatex -interaction=nonstopmode -halt-on-error class2PT26.tex
% Lecture topics follow PT26w/lectureNote.tex, sections 2 and 3.
\documentclass[a4paper,10pt,pdflatex,ja=standard,jafont=ipaex-type1,
  magstyle=real,scale=1]{bxjsarticle}
\usepackage{lmodern}
\usepackage{amsmath,amssymb,mathtools}
\usepackage{geometry}
\geometry{reset,left=22mm,right=22mm,top=20mm,bottom=17mm,
  headheight=13pt,headsep=6mm,footskip=10mm}
\usepackage{xcolor}
\usepackage{enumitem}
\usepackage{array,booktabs,tabularx}
\usepackage{fancyhdr}
\usepackage{hyperref}

\definecolor{ink}{HTML}{193A5C}
\definecolor{accent}{HTML}{237D82}
\definecolor{soft}{HTML}{EAF3F3}
\definecolor{muted}{HTML}{52616B}
\newcommand{\coursename}{確率論と統計学 I}
\newcommand{\lecturenumber}{2}
\hypersetup{unicode=true,colorlinks=true,urlcolor=accent,linkcolor=ink,
  pdfauthor={中島秀太},pdftitle={第2・3回：確率変数と分布、期待値と積分}}
\pagestyle{fancy}
\fancyhf{}
\fancyhead[L]{\small\color{ink}\coursename{}・授業要約と演習}
\fancyhead[R]{\small\color{ink}第\lecturenumber 回}
\fancyfoot[C]{\small\color{ink}\thepage/2}
\renewcommand{\headrulewidth}{0.3pt}
\setlength{\parindent}{0pt}
\setlength{\parskip}{2pt}
\setlength{\abovedisplayskip}{4pt}
\setlength{\belowdisplayskip}{4pt}
\setlength{\abovedisplayshortskip}{3pt}
\setlength{\belowdisplayshortskip}{4pt}
\setlist[itemize]{leftmargin=1.8em,itemsep=1pt,topsep=2pt,parsep=0pt}
\setlist[enumerate]{leftmargin=2em,itemsep=3pt,topsep=3pt,parsep=0pt}
\newcommand{\topic}[1]{\par\vspace{5pt}
  {\large\sffamily\mdseries\color{accent}#1}\par\vspace{2pt}}
\newcommand{\handouttitle}[1]{%
  \begin{center}
    {\Large\color{ink}#1}\par\vspace{3mm}
    {\color{accent}\rule{0.88\linewidth}{0.7pt}}
  \end{center}\vspace{-2mm}}
\newcommand{\keybox}[1]{\par\vspace{3pt}\begingroup
  \setlength{\fboxsep}{5pt}
  \colorbox{soft}{\parbox{\dimexpr\linewidth-2\fboxsep\relax}{#1}}%
  \endgroup\par\vspace{2pt}}
\newcommand{\examplelabel}[1]{{\sffamily\mdseries\color{ink}#1}}
\newcommand{\F}{\mathcal F}
\newcommand{\R}{\mathbb R}
\newcommand{\E}{\mathbb E}

\begin{document}
\fontsize{10pt}{13.5pt}\selectfont
\handouttitle{第2回\quad 確率変数と確率分布}

\topic{確率変数と可測性}
$(\Omega,\F,P)$ を確率空間、$\mathcal B(\R)$ を実数の開集合から生成される
最小の $\sigma$-加法族（ボレル集合族）とする。写像 $X:\Omega\to\R$ が
確率変数であるとは、次の可測性を満たすことである。
$$
  X^{-1}(B):=\{\omega\in\Omega:X(\omega)\in B\}\in\F
  \qquad\bigl(B\in\mathcal B(\R)\bigr).
$$
これは、観測値 $X$ に関する事象の確率を測れるという条件である。
可測性は $\{X\le x\}\in\F$ がすべての $x\in\R$ について成り立つことと同値である。

\topic{確率分布：観測値の空間上の確率測度}
$X$ の分布（押し出し測度）$\mu_X$ を、$\mathcal B(\R)$ 上で
$$
  \mu_X(B)=P(X\in B):=P\bigl(X^{-1}(B)\bigr)
  \qquad\bigl(B\in\mathcal B(\R)\bigr)
$$
と定める。離散確率空間では、各値 $x\in\R$ の確率は
$$
  P(X=x)=\sum_{\omega\in\Omega:\,X(\omega)=x}p_\omega.
$$
\keybox{確率変数 $X$ は写像、分布 $\mu_X$ は測度である。同じ分布をもつ確率変数でも、
標本空間や他の確率変数との関係は異なり得る。}

\topic{分布関数}
$F_X(x):=P(X\le x)=\mu_X(({-\infty},x])$ と定める。次の性質をもつ。
$$
  \begin{aligned}
    x\le y&\ \Longrightarrow\ F_X(x)\le F_X(y),
    &\lim_{h\downarrow0}F_X(x+h)&=F_X(x),\\
    \lim_{x\to-\infty}F_X(x)&=0,
    &\lim_{x\to+\infty}F_X(x)&=1.
  \end{aligned}
$$
特に $a<b$ に対して $P(a<X\le b)=F_X(b)-F_X(a)$ である。
分布関数は分布 $\mu_X$ を一意に定める。

\examplelabel{例：公平なコインを独立に2回投げ、表の個数を $X$ とする。}
$$
  X(H,H)=2,\quad X(H,T)=X(T,H)=1,\quad X(T,T)=0,
$$
$$
  P(X=0)=\frac14,\qquad P(X=1)=\frac12,\qquad P(X=2)=\frac14.
$$
したがって $F_X(0)=1/4$、$F_X(1)=3/4$、$F_X(2)=1$ となる。

\topic{演習問題}
\begin{enumerate}[label=\arabic*.]
  \item 公平なサイコロを1回投げ、出た目を $X$ とする。
  $X$ の分布と、すべての $x\in\R$ に対する分布関数 $F_X(x)$ を求めよ。
  \item $\Omega=[0,1]$、$\F=\mathcal B([0,1])$、$P=\lambda$（ルベーグ測度）とする。
  $X(\omega)=\omega$、$Y(\omega)=\omega+1$ に対し、$S=X+Y$ の分布関数を求めよ。
\end{enumerate}

\topic{レポート問題}
\begin{enumerate}[label=\arabic*.]
  \item 演習第2問と同じ確率空間と $X,Y$ に対し、$T=XY$ の分布関数を求めよ。
  $X,Y$ は同じ標本 $\omega$ の関数であることに注意せよ。
  \item 集合 $\{X\le x+1/n\}$ の列を用いて、$F_X$ の右連続性を示せよ。
  $A_n\downarrow A$ なら $P(A_n)\downarrow P(A)$ となる性質を用いてよい。
  \item $F_X(x-):=\lim_{y\uparrow x}F_X(y)$ とおく。
  $P(X=x)=F_X(x)-F_X(x-)$ を示し、分布関数の跳びの意味を説明せよ。
\end{enumerate}

\newpage
\renewcommand{\lecturenumber}{3}
\handouttitle{第3回\quad 期待値とルベーグ積分}

\topic{有限和から非負確率変数の積分へ}
以下、$(\Omega,\F,P)$ を確率空間とする。$A\in\F$ の指示関数 $\chi_A$ は
$A$ 上で1、それ以外で0をとる。非負単関数 $Z=\sum_{j=1}^m a_j\chi_{A_j}$
（$a_j\ge0$、$A_j\in\F$）について
$$
  \E[Z]:=\sum_{j=1}^m a_jP(A_j)
$$
と定める。この値は単関数の表示によらない。非負確率変数 $X$ について
$$
  \E[X]:=\int_\Omega X\,dP
  =\sup\{\E[Z]:0\le Z\le X,\ Z\text{ は非負単関数}\}\in[0,\infty]
$$
と定義する。下からの単関数近似 $0\le Z_n\uparrow X$ により積分を構成する。

\topic{可積分性と期待値の基本性質}
$X^+:=\max\{X,0\}$、$X^-:=\max\{-X,0\}$ とおく。
$$
  X=X^+-X^-,\qquad |X|=X^++X^-,\qquad
  \E[|X|]=\E[X^+]+\E[X^-].
$$
$\E[|X|]<\infty$ のとき $X$ は可積分であり、$\E[X]:=\E[X^+]-\E[X^-]$ と定める。
可積分な $X,Y$ と $a,b\in\R$ に対して
$$
  \E[aX+bY]=a\E[X]+b\E[Y],\qquad
  X\le Y\ \Longrightarrow\ \E[X]\le\E[Y],\qquad
  |\E[X]|\le\E[|X|].
$$

\topic{分布から期待値を計算する}
$\mu_X$ を $X$ の分布とする。ボレル可測関数 $g:\R\to\R$ について、
$g\ge0$ または $\E[|g(X)|]<\infty$ なら
$$
  \E[g(X)]=\int_\R g(x)\,\mu_X(dx).
$$
離散分布 $P(X=x_k)=q_k$ なら $\E[X]=\sum_k x_kq_k$。
密度 $f$、すなわち $\mu_X(B)=\int_B f(x)\,dx$ をもつなら
$\E[X]=\int_\R xf(x)\,dx$ である（いずれも非負または可積分の場合）。
\keybox{期待値は実際に出る値とは限らない。表なら100円、裏なら0円を得る公平なゲームでは、
期待値は50円だが、1回の結果は0円か100円である。}

\topic{極限と期待値の交換}
\examplelabel{単調収束定理：}
$0\le X_n\uparrow X$ がほとんど確実に成り立つなら $\E[X_n]\uparrow\E[X]$。
値 $+\infty$ も許される。

\examplelabel{優収束定理：}
$X_n\to X$ がほとんど確実に成り立ち、すべての $n$ について
$|X_n|\le Y$ がほとんど確実に成り立つ可積分な $Y\ge0$ があれば、
$X$ も可積分で $\E[X_n]\to\E[X]$。
「ほとんど確実に」とは、確率1の事象上で成り立つことである。

\topic{演習問題}
\begin{enumerate}[label=\arabic*.]
  \item 重み付きサイコロ $P(\{k\})=k/21$（$k=1,\ldots,6$）で、出た目 $X$ の期待値を求めよ。
  \item $\{1,2,3,4\}$ の24個の置換を等確率で選ぶ。不動点数
  $N(\sigma)=\#\{i:\sigma(i)=i\}$ を指示関数の和で表し、$\E[N]$ を求めよ。
\end{enumerate}

\topic{レポート問題}
以下、$\Omega=(0,1)$、$\F=\mathcal B((0,1))$、$P=\lambda$ とする。
\begin{enumerate}[label=\arabic*.]
  \item $X_n(\omega)=n\chi_{(0,1/n)}(\omega)$ とする。
  $X_n\to0$ と $\E[X_n]=1$ を示し、優収束定理を適用できない理由を述べよ。
  \item $X(\omega)=\omega^{-1/2}$、$X_n=\min\{X,n\}$ とする。
  $\E[X_n]$ を計算し、単調収束定理を用いて $\E[X]$ を求めよ。
\end{enumerate}
\end{document}
